Gestire un ambiente multi-cloud è una tema sempre più critico in letteratura. Dai dati dell'Osservatorio Cloud Transformation 2019 \cite{osservatorio} della School of Management del Politecnico di Milano, emerge che il 68\% delle imprese italiane costruisce un ambiente multicloud utilizzando in modo sistematico più di un Cloud Provider. Tuttavia, ciò che emerge è che solo il 24\% dichiara di riuscire a gestire i servizi di diversi provider in modo ottimale. Gli ostacoli al raggiungimento di una piena sinergia sono:
\begin{itemize}
    \item strumenti diversi per i flussi di lavoro e la gestione;
    \item mancanza di sicurezza unificata;
    \item notevole aumento della complessità del sistema;
    \item difficoltà nella sincronizzazione e nella condivisione dei dati.
\end{itemize}
L'aumento della complessità di gestione di un'ambiente multi-cloud ha portato alla nascita di varie tecnologie  che permettono di automatizzare molti di questi processi. Oggigiorno sono molte le soluzioni presenti sul mercato, tra le più importanti vi sono le soluzioni offerte da IBM \cite{ibm} e VMware \cite{vmware}. La soluzione multi-cloud di IBM permette di facilitare, orchestrare e ottimizare gli ambienti con più cloud sia pubblici che privati che con infrastrutture locali, contenere al massimo i costi, pur incrementato il vantaggio di business. VMware offre le funzionalità operative multi-cloud necessarie per utilizzare qualsiasi cloud per ogni app, incluse macchine virtuali e container su Private Cloud e Hybrid Cloud basati su VMware e carichi di lavoro Public Cloud nativi su AWS, Azure, Google, Oracle e altri. 